% ESTADO DEL ARTE Y MARCO TEÓRICO

\section{Estado del Arte y Marco teórico}

El estado del arte es el conocimiento más actualizado que existe para dar una solución al problema planteado en una investigación. Se explican las diferentes alternativas tecnológicas, así como las ventajas y desventajas en su implantación en el trabajo. Es probable que al comenzar un proyecto se requiera de un análisis de las tecnologías disponibles, así como metodologías o procesos pueden estar implicados para la consecución de los objetivos. Esta tarea de investigación supone un ejercicio de comprensión de la dificultad de las tareas y descripción de las alternativas disponibles en otros trabajos.

El marco teórico es el conocimiento mínimo que resulta indispensable y necesario para poder comprender el problema y las características del mismo.

No es necesario mantener el título “Estado del Arte”. Se puede sustituir por un titulo relacionado con las bases teóricas del proyecto.

Añadir todas las secciones que sean necesarias para presentar una estructura coherente.

Las imágenes deben estar referidas en el texto para que cualquier descripción acompañada de una figura pueda entenderse en el contexto. Ejemplo: Se habilita un servidor con Jupyter como se puede ver en la Figura 1, en el cual se puede disponer de diferentes kernels de programación (Python, R, Julia…).

\begin{figure}[h!]
  \centering
  \includegraphics[width=13.3cm]{fig1.png}
\caption[Arquitectura Jupyter Cliente - Servidor]{Arquitectura Jupyter Cliente - Servidor. Fuente: https://www.paradigmadigital.com/dev/jupyter-data-science-aplicada/}
  \label{fig:arquitectura}
\end{figure}

Utilizar el estilo “Descripción” para insertar las figuras.

Está permitido el uso de viñetado:

\begin{itemize}
  \item Esta sería la forma
  \item Segundo
  \begin{itemize}
    \item En caso de tener un subapartado
  \end{itemize}
\end{itemize}

Una numeración que se quiera hacer no debe influir en la numeración de los apartados generales:

\begin{enumerate}
  \item Item 1
  \item Item 2
\end{enumerate}
